%% kapa: exam=1 round=21 year=2010 period=1 form=A subject=civil_law subject_ko=민법 seq=19 number=19
%% source: Q-Net 감정평가사 기출문제 (cst00309) artlSeq=5024181 file=extracted/1008/5024181/2010년 제21회 감정평가사 1차시험 기출문제/1교시.hwp
%% conversion: kordoc 4.24.51 -> tools/gen1.py (auto); review=proofread-sonnet; flags=none
\begin{question}{19}
대리에 관한 설명으로 옳은 것을 모두 고른 것은?
\begin{framedtext}
ㄱ. 대리권의 범위가 정해지지 않은 대리인은 대리의 목적인 물건이나 권리의 성질을 변화시키지 않는 범위에서 개량행위를 할 수 있다. \par
ㄴ. 임의대리인은 본인의 승낙이 있거나 부득이한 사유가 있는 때가 아니면 복대리인을 선임할 수 없다. \par
ㄷ. 법정대리인은 원칙적으로 복대리인을 선임할 수 없으며, 법원의 허가를 얻어 예외적으로 복대리인을 선임할 수 있다. \par
ㄹ. 임의대리의 종료사유는 본인의 사망, 대리인의 사망이나 금치산 또는 파산, 임의대리의 원인된 법률관계가 종료된 경우로 한정된다. \par
ㅁ. 무권대리의 상대방은 상당한 기간을 정하여 본인에게 추인여부의 확답을 최고할 수 있으며, 그 기간 내에 본인이 확답을 발하지 않으면 추인을 거절한 것으로 본다. \par
\end{framedtext}
\begin{choices}
  \item ㄱ, ㄴ
  \item ㄴ, ㄷ
  \item ㄱ, ㄴ, ㅁ
  \item ㄷ, ㄹ, ㅁ
  \item ㄱ, ㄷ, ㄹ, ㅁ
\end{choices}
\end{question}
\answer{3}
